% International Journal of Performance Analysis in Sport (Taylor & Francis) version.
% Do not compile this file directly: compile anonymous.tex (double-anonymous review copy) or authors.tex (copy with
% author details). Both set \ifanonymous and then \input this file. Interact layout + APA references, as for SMF
% (format-free submission). Every number comes from docs/results/*.md (rows 1-26 of docs/PAPER_V2_NOTES.md).
\documentclass[]{interact}

\usepackage{graphicx}
\usepackage{epstopdf}
\usepackage[caption=false]{subfig}
\usepackage[natbibapa,nodoi]{apacite}
\setlength\bibhang{12pt}
\renewcommand\bibliographytypesize{\fontsize{10}{12}\selectfont}
\usepackage{amsmath}
\usepackage{booktabs}
\usepackage{xcolor}
\usepackage{lineno}
\usepackage{hyperref}
\ifanonymous\hypersetup{pdfauthor={},pdfsubject={},pdfkeywords={}}\fi
% Word counts, filled in by paper/journal/ijpas/make_ijpas.py (counts.tex); fallback if the file is missing.
\IfFileExists{counts.tex}{% generated by make_ijpas.py
\newcommand{\wordcount}{5,519}
\newcommand{\abstractcount}{186}
}{\newcommand{\wordcount}{?}\newcommand{\abstractcount}{?}}

\begin{document}

\articletype{RESEARCH ARTICLE}

\title{How much can you trust broadcast pressing data? Visibility bias in team pressing metrics, its drivers, and a
correction validated on real broadcast video}

\ifanonymous
\author{}
\else
\author{
\name{Varun Inamdar\thanks{CONTACT Varun Inamdar. Email: vninamdar03@gmail.com}}
\affil{Vishwakarma University, Pune, India}
}
\fi

\maketitle

\noindent\textbf{Word count:} \wordcount{} words (main text, tables, figure and table captions and references; abstract
\abstractcount{} words).

\begin{abstract}
Broadcast-derived tracking data is widely available, but the main camera shows only about 14 of the 22 players at a
time. In this simulation-based measurement study, we quantified how this biases pressing metrics and how to correct it.
Using optical tracking of 10 professional matches (two providers, 14 teams), we hid every player that a camera fitted
to 20 real broadcasts would not show, and recomputed pressing. Broadcast tracking undercounted a team's active
pressers by 22.6\% (95\% CI 19.8--25.2) and 20.2\% (17.1--23.2) in the two datasets, and every team was undercounted.
The number of visible players almost fully explained the undercount ($r = 0.98$; about 4 percentage points per
player), whereas pressure on the ball carrier was nearly unbiased unless the camera lagged. A visible-player
correction reduced the mean absolute error from 21.9\% to 3.0\% for unseen cameras and matches. On real video, an
open pipeline counted visible players within half a player on two independent datasets, and filtering out close-ups
kept the correction within 1 percentage point on eight full broadcasts. Team-wide pressing metrics from broadcast
data should be corrected, or compared only between similarly framed broadcasts.
\end{abstract}

\begin{keywords}
football; broadcast tracking; pressing; measurement bias; performance analysis; computer vision
\end{keywords}

% =============================================================================================
\section{Introduction}

Pressing, how intensely a team without the ball closes down the opponent, is one of the most discussed tactical
concepts in modern football. Quantifying it requires the positions of all players. Full optical tracking provides
them, but most clubs, analysts and researchers cannot access it. Tracking data extracted from broadcast video is a much
cheaper alternative. It is increasingly offered commercially and studied academically
\citep{somers2024gsr,smf2025broadcast}. Its limitation is structural: the broadcast camera follows the action, so
players far from the ball are often off screen.

For spatial metrics such as pitch control, this partial view causes large errors in the hidden areas of the pitch
\citep{offscreen2026}, and several methods impute off-screen players
\citep{omidshafiei2022graphimputer,rsos2025continuous,offscreen2026}. Comparing broadcast-derived with official
tracking, \citet{smf2025broadcast} found that less than half of the players' activity is visible on TV broadcasts.
They concluded that some tactical analyses remain reliable, while physical metrics are more affected. For pressing
metrics, the size of the bias, what drives it and how certain we can be about it have not been measured. Pressing
involves the defenders near the ball and also those covering passing lanes further away. It is therefore not obvious
whether a ball-following camera biases it at all.

This study addressed five questions. (1) How large is the visibility bias in pressing metrics, and how certain is it?
(2) Is it driven by particular teams, matches or stadium cameras? (3) Which measurable properties of a broadcast drive
it? (4) Can it be corrected, also for broadcasts filmed in another league? (5) Does the correction work on real
broadcast video, where its input must be measured by computer vision and replays and close-ups interrupt the main
camera?

% =============================================================================================
\section{Methods}

\subsection{Data}

\textbf{Ground-truth tracking.} The IDSSE open dataset \citep{bassek2025idsse} (CC-BY 4.0) contains optical tracking
of all 22 players and the ball at 25\,Hz for 7 matches of the 2022/23 Bundesliga and 2.~Bundesliga, with possession
information (552{,}294 ball-in-play frames). As a second, independent provider we used the 3 Metrica Sports sample
games \citep{metrica_sample} (25\,Hz, all players). Their possession and ball-in-play status were reconstructed from
the accompanying event files. Together the ground truth covers 10 matches and 14 teams. One team, Fortuna
D\"usseldorf, appears in 5 of the 7 IDSSE matches; Section~\ref{sec:teams} tests whether it drives the results. All
tracking data were loaded with kloppy \citep{kloppy}.

\textbf{Real broadcast footprints.} The SkillCorner open data \citep{skillcorner_opendata} (MIT licence) contains
broadcast tracking of 20 A-League 2024/25 matches at 10\,Hz. For every frame it gives the four corners of the camera
image projected onto the pitch, and for every player whether they were detected on screen. These data were used only
to fit and validate the camera model (375{,}357 frames).

\textbf{Broadcast video.} SoccerNet-GSR \citep{somers2024gsr} (GPL-3.0) contains 30-second main-camera clips from the
Swiss Super League (25\,Hz, 1920$\times$1080). For every frame it labels each person's box, role, team, jersey number
and pitch position. We trained on the 57 training clips, tuned on half of the 58 validation clips and report on the
other half and on the 49 test clips. As a fully independent test of the visible-player count we used the 30 labelled
football clips of SportsMOT \citep{cui2023sportsmot} (CC BY-NC 4.0; 20{,}171 frames, 1280$\times$720). They come
from other broadcasts, and only players on the pitch are labelled. The labelled clips of SoccerNet-Tracking were not
used: all 106 are also SoccerNet-GSR clips.

\textbf{Broadcast composition.} The SoccerNet-v2 camera-shot labels \citep{deliege2021soccernetv2} cover 500 full
broadcasts (2014--2017; Premier League, La Liga, Serie A, Bundesliga, Ligue 1 and Champions League). Many are labelled
only around selected moments. We therefore used the 203 broadcasts with at least 300 labelled camera changes to check
how labels are timed, and the 196 of these without an unlabelled gap longer than 3 minutes (297 hours) to measure how
much of a broadcast is not the live main camera. No video frames are reproduced in this article.

\subsection{Pressing metrics}
For each frame with a ball carrier we computed the matrix $P$ of interception probabilities between defenders and
attackers \citep{bekkers2025pressing}, using \texttt{unravelsports} \citep{unravelsports} with a speed threshold of
2\,m/s. We report two summaries per frame. \emph{Pressure on the ball carrier}, $1 - \prod_i (1 - P_{i,c})$ over
defenders $i$ for carrier $c$, is a local metric. The number of \emph{active pressers}, defenders whose largest
probability exceeds 0.5, is a team-wide metric. Frames without a ball carrier in the possessing team were dropped:
0.06--0.2\% per IDSSE match and 6.8\% for Metrica, whose possession comes from hand-annotated events.

\subsection{A broadcast camera fitted to real footage}
The main camera was modelled as a fixed point $(c_x, c_y)$ behind the near touchline. A player at bearing $\theta$
from the camera is visible if $|\theta - \phi_t| \le \alpha_t/2$ and their depth along the viewing direction exceeds
the bottom image edge $\nu_t$:
\begin{align}
\phi_t &= \mathrm{EMA}_{\lambda}\!\left(g_b\,\theta^{\text{ball}}_t + g_p\,\bar\theta^{\text{players}}_t + o\right),\\
\alpha_t &= \mathrm{clip}\!\left(a_0 + a_1 d^{\text{ball}}_t + a_2\, s^{\text{players}}_t\right),\qquad
\nu_t = n_0 + n_1 d^{\text{ball}}_t.
\end{align}
Here $\phi$ is the pan and $\alpha$ the horizontal field of view. $d^{\text{ball}}$ is the ball's distance from the
camera, $\bar\theta^{\text{players}}$ and $s^{\text{players}}$ are the mean and standard deviation of the players'
bearings, and $\lambda$ is an exponential lag.

In every SkillCorner frame, the left and right image edges are rays from the camera's ground position; their
intersection locates the camera (median per match). Pan and field of view follow from the bearings of the far image
corners, and the bottom edge from the near corners. The coefficients were fitted by least squares and $\lambda$ by
grid search (0--1\,s). One shared camera was fitted on 10 matches and tested on the other 10 against the per-player
detection flags. To capture differences between stadiums we also fitted one camera per match (20 stadium cameras).
The commonly assumed baseline is a full-width window of fixed length centred on the ball.

\subsection{Simulating the broadcast view and estimating the bias}
Given full tracking and a camera, we removed every player outside the camera's view and recomputed the metrics on the
remaining players. The ball and the ball carrier were always kept, and a frame with no visible defender counted as
zero pressure. The bias of a metric is the relative difference between a team's out-of-possession mean with and
without hiding. Each stadium camera was applied to each ground-truth match, giving 400 team-broadcasts
(20 cameras $\times$ 10 matches $\times$ 2 teams).

\subsection{Correction}
Because the bias of active pressers was nearly linear in the number of visible players (Section~\ref{sec:drivers}),
we corrected a broadcast's value by fitting $\text{bias} = k \cdot \text{visible} + b$ across broadcasts and dividing
by $1 + \text{bias}$. The correction was evaluated strictly. For every team-broadcast, the line was fitted without its
stadium camera and without its match. In a stricter variant, every match of its pressing team was also left out.

\subsection{Robustness and generalisation}
\label{sec:robust-methods}
\textbf{Teams.} We report the bias per team, a team-balanced mean (every team weighs the same) with an interval from
resampling teams, and leave-one-team-out estimates. We also estimated the extra bias of the most frequent team at
equal visibility, from the least-squares model $\text{bias} = a + b\cdot\text{visible} + c\cdot[\text{team}]$ with a
match-level bootstrap.

\textbf{A second league.} SoccerNet-GSR gives no camera parameters. However, each person's labelled pitch position
is an exact projection of the bottom of their box. A homography fitted to these positions per frame therefore maps the
image corners onto the pitch, giving the same footprints as SkillCorner. We fitted the camera model in this way to 6
Swiss Super League broadcasts. Off-screen players are unknown in this data, so the pan followed the ball only. The
Swiss cameras were applied to all 10 ground-truth matches (120 team-broadcasts). These broadcasts were corrected with
the line fitted on the A-League cameras, leaving out each broadcast's match.

\subsection{Measuring visible players on real video}
The correction needs one number per broadcast: the mean number of visible players per frame (both teams and the
goalkeepers, no referees). We measured it with an open pipeline.
\begin{itemize}
\item \emph{Detection.} An RF-DETR Medium detector \citep{rfdetr} fine-tuned on SoccerNet-GSR (95.8\% recall, 95.5\%
precision for people on validation).
\item \emph{Tracking.} OC-SORT \citep{cao2023ocsort} with camera-motion compensation \citep{trackers}.
\item \emph{Teams and roles.} Teams from per-clip $k$-means on shirt colour. Tracks that fit neither team are
goalkeepers inside a penalty area and referees elsewhere.
\item \emph{Pitch calibration.} A ResNet-34 \citep{he2016resnet} pitch-landmark network with a RANSAC homography.
\item \emph{Jersey numbers.} A fine-tuned scene-text recogniser (PARSeq; \citealp{bautista2022parseq}).
\end{itemize}
The count of visible players uses the tracked people whose role is player or goalkeeper. It does not need a
homography. A linear calibration ($\text{labels} \approx a\cdot\text{count} + b$) was fitted on the SoccerNet-GSR
tuning clips only and applied unchanged to all other data.

The counting error was propagated into the correction. For each team-broadcast we added count errors drawn from the
held-out clips to its true visible-player mean (500 draws). A broadcast is a whole match, about 180 clips, so we
averaged 180 clip errors per draw. As a worst case we also used the error of a single clip.

\textbf{Replays and close-ups.} Non-main-camera frames must not enter the count. A main-camera filter removes shots
with little grass ($<0.2$ of the frame) or very large people (median box over 300\,px). Shot boundaries come from
colour-histogram cuts, and each shot is kept or dropped as a whole. From the SoccerNet-v2 labels we measured the
share of broadcast time per camera type. On the 203 broadcasts we verified that a label describes the shot ending at
its timestamp: at shots on target and goals, 99\% of labels read this way are the live main camera. We then derived how
many non-main frames the filter may let through before the correction moves by more than 1 percentage point.
Finally, we tested the filter on eight full broadcasts from six competitions (SoccerNet raw videos, 720p at 2 frames
per second, 86{,}400 frames), scoring every frame against its camera label. The 85{,}117 frames inside a labelled
shot were scored; the rest lie before the first or after the last labelled camera change of a half. People were
detected on every second frame. We report the share of each camera type removed and the shift of the mean people count on the kept frames
relative to the live main-camera frames, converted to percentage points with the slope of 4.1. We then tuned the
thresholds (grass share, box height, shot voting, and whether frames without detected people pass) leaving one
broadcast out: the setting with the smallest mean shift on seven broadcasts, removing at most 5\% of their live
main-camera frames, was scored on the eighth. Only aggregate statistics were stored; no video frames were retained.

\subsection{Statistics}
Frames 40\,ms apart are strongly dependent, so we used a block bootstrap with 60\,s blocks within each team-match.
Across stadium cameras, a three-level bootstrap resampled cameras, matches (both teams together) and blocks
(2{,}000 replicates, 95\% percentile intervals). For the stadium-camera runs, metrics were evaluated every 0.2\,s
while the camera followed every frame. This changed the bias by less than 0.1 percentage points (checked on one full
match). Counting errors on video are reported with intervals from resampling whole clips.

% =============================================================================================
\section{Results}

\subsection{The fitted broadcast camera}

\begin{figure}[t]
\centering
\includegraphics[width=\textwidth]{figures/fig1_cameras.pdf}
\caption{The same frame seen by the commonly assumed rectangle camera (left) and the camera fitted to real broadcasts
(right). The fitted camera stands behind the near touchline: its footprint widens towards the far side and cuts off
players close to the camera.}
\label{fig:cameras}
\end{figure}

The fitted camera stood on the halfway line, 57--97\,m from the pitch centre depending on the stadium. Its pan
followed the ball and the players' centre with similar weights ($g_b = 0.57$, $g_p = 0.56$; $R^2 = 0.97$). The shared
camera showed no measurable lag. The median field of view was 27$^\circ$ (19--40$^\circ$); it narrowed when play was
far away and widened when players were spread out (Figure~\ref{fig:cameras}). On 10 held-out matches the fitted camera
predicted which players were on screen with 87.1\% agreement, against 82.9\% for a 40\,m rectangle and a 96.8\% upper
bound from the real footprint. It beat the rectangle on every held-out match and reproduced the mean number of visible
players exactly (14.65; Table~\ref{tab:validation}). A camera following only the ball underestimated visibility
(13.30 players), showing that camera operators frame the players, not only the ball.

\begin{table}[t]
\centering
\caption{Which players are on screen? Agreement of each camera model with SkillCorner's per-player detection flags on
10 held-out matches.}
\label{tab:validation}
\begin{tabular}{lcc}
\toprule
Camera model & Agreement & Visible players per frame \\
\midrule
Real footprint (upper bound) & 96.8\% & -- \\
\textbf{Fitted camera} & \textbf{87.1\%} & \textbf{14.65} \\
Rectangle, 40\,m & 82.9\% & 15.76 \\
Rectangle, 32\,m with assumed lag, perspective and zoom & 82.2\% & 15.29 \\
Fitted camera following the ball only & 82.0\% & 13.30 \\
\midrule
Real broadcast & -- & 14.65 \\
\bottomrule
\end{tabular}
\end{table}

\subsection{Visibility bias in pressing metrics}

\begin{table}[t]
\centering
\caption{Relative bias of pressing metrics under different cameras (negative = undercount), with 95\% bootstrap
intervals. The stadium-camera interval includes camera, match and within-match variation.}
\label{tab:bias}
\small
\begin{tabular}{llccc}
\toprule
Data & Camera & Visible & Pressure on carrier & Active pressers \\
\midrule
IDSSE (7) & Rectangle 40\,m & 14.5 & $-0.1\%$ [$-0.2$, $-0.1$] & $-20.3\%$ [$-22.7$, $-17.9$] \\
IDSSE (7) & Rectangle 40\,m, assumed 0.5\,s lag & 14.3 & $-4.1\%$ [$-4.6$, $-3.7$] & $-24.1\%$ [$-26.5$, $-21.9$] \\
IDSSE (7) & Fitted shared camera & 13.8 & $-0.5\%$ [$-0.7$, $-0.4$] & $-21.4\%$ [$-23.8$, $-19.1$] \\
IDSSE (7) & \textbf{20 stadium cameras} & 12.1--16.0 & $\mathbf{-1.7\%}$ [$-2.2$, $-1.2$] & $\mathbf{-22.6\%}$ [$-25.2$, $-19.8$] \\
Metrica (3) & \textbf{20 stadium cameras} & 12.7--16.6 & $\mathbf{-1.8\%}$ [$-2.3$, $-1.3$] & $\mathbf{-20.2\%}$ [$-23.2$, $-17.1$] \\
\bottomrule
\end{tabular}
\end{table}

With every camera model, broadcast tracking undercounted a team's active pressers by about one fifth
(Table~\ref{tab:bias}). Averaged over the 20 stadium cameras, the undercount was 22.6\% for IDSSE and 20.2\% for
Metrica. Pressure on the ball carrier was much less affected ($-0.5\%$ to $-1.8\%$). The stadium camera alone
contributed as much uncertainty as the sample of matches (IDSSE, active pressers: camera only [$-24.3$, $-20.6$];
matches and blocks only [$-24.4$, $-20.6$]). The same camera showed about one more player on Metrica than on IDSSE.
At equal visibility the two providers agreed: at 14.5 visible players, the undercount of active pressers was 19.6\% in
both datasets.

\subsection{What drives the bias}
\label{sec:drivers}

\begin{figure}[t]
\centering
\begin{minipage}{0.49\textwidth}\centering
\includegraphics[width=\linewidth]{figures/fig2_visibility.pdf}
\end{minipage}\hfill
\begin{minipage}{0.49\textwidth}\centering
\includegraphics[width=\linewidth]{figures/fig3_lag.pdf}
\end{minipage}
\caption{One point per stadium camera and dataset. Left: the undercount of active pressers shrinks by about
4 percentage points per additional visible player ($r = 0.98$). Right: the bias in pressure on the ball carrier grows
with the camera's fitted lag; $n$ = number of stadium cameras per lag value.}
\label{fig:drivers}
\end{figure}

Across stadium cameras, the bias was almost fully explained by two properties of the broadcast
(Figure~\ref{fig:drivers}). For active pressers it was the \textbf{number of visible players}. Broadcasts that framed
the whole group (15--16.6 players visible) missed 11--14\%, and tight framing (about 12 visible) missed up to 29\%
($r = 0.98$ in both datasets; 3.9 and 4.2 percentage points per player). For pressure on the ball carrier it was
\textbf{camera lag}. The 11 stadium cameras without measurable lag gave $-0.9\%$; 0.1\,s of lag gave $-1.7\%$, 0.2\,s
$-3.0\%$ and 0.3\,s $-3.5\%$: about $-1.0$ percentage point per 0.1\,s of lag (95\% CI $-1.2$ to $-0.9$, from
resampling the 20 cameras; $r = -0.97$). When the camera trails the ball, defenders closing down the
carrier enter the picture late.

\subsection{Correcting the bias}
\label{sec:correction}
Over the 400 team-broadcasts, the fitted correction line was
$\text{bias} = 4.12 \cdot \text{visible} - 79.3$ (in \%; slope 95\% CI 3.68 to 4.52 and intercept $-85.4$ to $-72.9$,
from resampling cameras and matches; $R^2 = 0.76$). With the stadium camera
and the match held out, the correction reduced the mean absolute error in active pressers from 21.9\% to
\textbf{3.0\%}, and 95\% of broadcasts were within $\pm 7.0\%$. With the pressing team also unseen, the error was
3.2\% (95th percentile 7.6\%). Per match, the corrected error ranged from 1.5\% to 4.9\%.

\subsection{Robustness to teams}
\label{sec:teams}
Every one of the 14 teams was undercounted, by 17.8--27.9\%. The team-balanced mean was $-22.1\%$
[$-23.8$, $-20.4$] for IDSSE and $-20.8\%$ [$-22.5$, $-19.0$] for Metrica. Leaving out all matches of any one IDSSE
team gave mean biases between $-19.5\%$ (without Fortuna D\"usseldorf, 2 matches left) and $-23.2\%$, with slopes of
3.9--4.2 percentage points per player. At equal visibility, Fortuna's extra bias was $-1.9$ percentage points
[$-5.3$, $+1.5$], so it did not differ from the other teams.

\subsection{A second league's cameras}
\label{sec:league}

\begin{figure}[t]
\centering
\includegraphics[width=\textwidth]{figures/fig5_generalisation.pdf}
\caption{(a) Bias in active pressers for cameras fitted to six Swiss Super League broadcasts (squares; mean over the
10 ground-truth matches) against the A-League cameras and their correction line. The shaded band is the visibility
range the line was fitted on. (b) How far the corrected value moves when the main-camera filter keeps a share of the
non-main-camera frames, for a typical broadcast (35\% non-main) with 14 visible players on main-camera frames.}
\label{fig:generalisation}
\end{figure}

The six Swiss cameras stood 51--66\,m from the pitch centre, with median fields of view of 24--34$^\circ$. Following
only the ball, they showed 9.8--12.8 players. That is fewer than the 11.5--16.0 players visible in the real Swiss
broadcasts, because the players' influence on the pan could not be fitted. Their bias in active pressers was
27--45\%, and it fell more steeply with visibility than for the A-League cameras (5.1 against 4.1 percentage points
per player; Figure~\ref{fig:generalisation}a). With the A-League line, the corrected mean absolute error was
\textbf{4.7\%} for broadcasts within the visibility range the line was fitted on ($\ge 11.7$ players). Below that
range it was 8.8\%, an under-correction that a quadratic fit did not remove. Fitting the line on both leagues
together brought the error to 4.6\% (95th percentile 11.5\%).

\subsection{From real broadcast video}

\begin{table}[t]
\centering
\caption{End-to-end game state on SoccerNet-GSR test (49 clips), scored with the official evaluator. GS-HOTA matches
a prediction to a labelled person only within 5\,m and with the same role, team and jersey number
\citep{luiten2021hota}. The position-only score ignores the attributes and uses our implementation of the same
similarity.}
\label{tab:gsr}
\begin{tabular}{lcc}
\toprule
Pipeline & GS-HOTA & Position only \\
\midrule
SoccerNet-GSR baseline \citep{somers2024gsr} & 22.3 & 42.7 \\
Ours, no jersey numbers & 24.0 & 62.3 \\
Ours, with jersey numbers & \textbf{39.8} & 62.3 \\
\bottomrule
\end{tabular}
\end{table}

\textbf{Pipeline.} Game-state reconstruction is not the aim of this study; Table~\ref{tab:gsr} only shows that the
pipeline is competitive (well above the public baseline). The output used here is the count of visible players.

\begin{table}[t]
\centering
\caption{Counting visible players on real video: mean error (pipeline minus labels, players per frame) with 95\%
intervals from resampling clips, error of 30-second clip means, and correlation of clip means. The calibration was
fitted on SoccerNet-GSR tuning clips only.}
\label{tab:count}
\small
\begin{tabular}{llccc}
\toprule
Data & Count & Mean error [95\% CI] & Clip MAE & Clip $r$ \\
\midrule
SoccerNet-GSR, held-out valid (29 clips) & raw & $-1.05$ [$-1.41$, $-0.75$] & 1.07 & 0.95 \\
SoccerNet-GSR, held-out valid (29 clips) & calibrated & $-0.32$ [$-0.68$, $-0.01$] & 0.67 & 0.95 \\
SportsMOT, independent (30 clips) & raw & $-0.26$ [$-0.54$, $-0.01$] & 0.70 & 0.90 \\
SportsMOT, independent (30 clips) & calibrated & $+0.53$ [$+0.26$, $+0.79$] & 0.78 & 0.90 \\
\bottomrule
\end{tabular}
\end{table}

\textbf{Counting visible players.} On SoccerNet-GSR the pipeline undercounted by about one player (small or occluded
players), and the calibration removed this bias (Table~\ref{tab:count}). On SportsMOT the raw count was already close,
so the calibration over-corrected; adding SportsMOT's own training clips to it barely helped ($+0.42$). The offset
therefore depends on the dataset, but the calibrated count was within about half a player on both datasets, and clip
means followed the labels closely ($r = 0.90$--$0.95$).

\begin{figure}[t]
\centering
\includegraphics[width=\textwidth]{figures/fig4_real_video.pdf}
\caption{(a) Visible players per 30-second clip: calibrated pipeline count against labels (SoccerNet-GSR, held-out
validation and test). (b) Error in active pressers after correction for the 400 team-broadcasts, with the true
visible-player mean, with the calibrated video count over a whole match or a single clip, and with the uncalibrated
count. Boxes: interquartile range; whiskers: 2.5--97.5\%.}
\label{fig:video}
\end{figure}

\textbf{The correction with the video count.} With the calibrated SoccerNet-GSR count over a whole match, the corrected
error equalled that with the true count (Figure~\ref{fig:video}b; mean absolute error 3.0\%). Without calibration it
overshot (mean error $+4.7\%$); from a single 30-second clip it was 4.6\%. With SportsMOT's count errors, the
whole-match error was 3.3--3.6\%.

\textbf{Broadcast composition and the main-camera filter.} In the 196 fully labelled broadcasts, the live main camera
made up 64\% of broadcast time (50--74\% per game; 57--67\% per competition), close-ups 29\% and main-camera replays
1.9\%. Main-camera replays show players as live play does; close-ups do not, and if all non-main frames entered the
count, the correction would move by 9--24 percentage points. To stay within 1 percentage point, the filter must remove
about 97\% of the non-main frames (Figure~\ref{fig:generalisation}b).

The default filter removed no frame of SoccerNet-GSR, which is all main camera, but on eight full broadcasts it kept
42\% of the non-main frames (Table~\ref{tab:filter}): the detector, trained on wide views, often found no one in a
close-up, and such frames passed the size test. The correction would move by $-7.1$ percentage points. Requiring at
least one detected person of at most 200\,px, a grass share of at least 0.15 and decisions per frame removed 93\% of
the close-ups and 1.7\% of the live main-camera frames; every leave-one-broadcast-out fold chose this setting, and the
held-out correction moved by 0.3 percentage points on average (at most 0.6). The frames that still passed were mostly
main-camera replays and wide views, whose counts match live play. On SoccerNet-GSR this filter removed 1.2\% of frames.

\begin{table}[t]
\centering
\caption{Main-camera filter on eight full broadcasts (85{,}117 labelled frames). Shift: change of the corrected
pressing value caused by the frames the filter keeps; held out: thresholds tuned on the other seven broadcasts.}
\label{tab:filter}
\begin{tabular}{lcc}
\toprule
 & Default filter & Broadcast filter \\
\midrule
Live main camera removed & 1.2\% & 1.7\% \\
Close-ups removed (27\% of frames) & 61\% & 93\% \\
Views behind the goal removed (3\%) & 36\% & 62\% \\
Main-camera replays removed (3\%) & 9\% & 10\% \\
Non-main frames kept & 42\% & 10\% \\
People count on kept frames vs live play & $-1.72$ & $+0.01$ \\
Shift, all broadcasts (pp) & $-7.1$ & $+0.03$ \\
Shift per broadcast, held out: mean $|\cdot|$ (max) (pp) & -- & 0.3 (0.6) \\
\bottomrule
\end{tabular}
\end{table}

% =============================================================================================
\section{Discussion}

Broadcast tracking undercounted team-wide pressing by about one fifth across 10 matches, 14 teams, two tracking
providers, 20 stadium cameras and a second league's cameras. The undercount was not a property of teams but of how
tightly the broadcast was framed, and carrier-centred pressing was biased only by camera lag. Both drivers can be
measured on the broadcast itself.

\textbf{Practical implications.} Analysts should treat team-wide pressing counts from broadcast tracking, such as the
number of players involved in a press, as systematically low. Comparing teams or matches filmed with different camera
styles can mislead: a tightly framed broadcast (about 12 visible players) undercounts by about 29\%, a wide one (about
16) by about 13\%. Carrier-centred metrics can be used with small error unless the camera lags (about 3\% at 0.2\,s).

\textbf{Applying the correction.} The correction needs only the mean number of visible players over the match. For
example, a broadcast showing 12.5 players on average has a predicted bias of $4.12 \times 12.5 - 79.3 = -27.8\%$, so an
observed 3.0 active pressers per frame corresponds to $3.0 / (1 - 0.278) \approx 4.2$. The count can come from an open
pipeline once close-ups are filtered out, and its uncertainty ($\pm 0.5$ players) adds less than one percentage point
over a whole match. The correction should be used within the fitted visibility range (about 12--17 players), below
which it under-corrects, and for match-level means rather than short excerpts.

\textbf{Relation to previous work.} Imputing off-screen players \citep{omidshafiei2022graphimputer,offscreen2026}
targets individual positions and needs a trained model. Our correction targets the aggregate metric directly and
needs a single number per broadcast. Our findings extend \citet{smf2025broadcast}: they concluded that tactical
analyses are less affected than physical ones, but a central tactical metric, team-wide pressing, is strongly
biased.

\textbf{Limitations.}
\begin{itemize}
\item \emph{The broadcast view is simulated} on ground-truth tracking. No open dataset contains broadcast video and
full tracking of the same match, so we validated the correction's input on real video, not the corrected pressing
value itself.
\item \emph{Camera fits.} The camera was fitted on A-League broadcasts. The Swiss cameras could be fitted only to the
ball and are therefore narrower than the real ones. SkillCorner's detection flags also reflect detector errors, which
caps agreement at about 97\%.
\item \emph{Ground truth.} It covers 10 matches, with one team in 5 of the 7 IDSSE matches, although
Section~\ref{sec:teams} shows that this team does not drive the result. Metrica possession is reconstructed from
hand-annotated events.
\item \emph{Count calibration.} Its offset depends on the dataset, so we report a count uncertainty of
$\pm 0.5$ players rather than a single calibration.
\item \emph{Broadcast era.} The labelled broadcasts are from 2014--2017, and current broadcasts may be composed
differently.
\item \emph{Replay filter.} It was tuned and tested on eight broadcasts from 2015--2017 with our detector at 2
frames per second. It keeps most main-camera replays, which is harmless for the count but would matter for metrics
that must not count an action twice. Other broadcasters, or a detector that finds people in close-ups, may need the
thresholds re-tuned with the same procedure.
\end{itemize}

% =============================================================================================
\section{Conclusion}

Broadcast tracking undercounts a team's active pressers by about one fifth, and by up to 29\% for tightly framed
broadcasts, while pressure on the ball carrier is nearly unbiased unless the camera lags. The undercount is almost
fully explained by the number of visible players. A correction based on that number reduces the error from 21.9\% to
3.0\% on unseen stadiums and matches, and to about 5\% on another league's cameras within the fitted visibility
range. On real broadcast video, the number of visible players can be measured to within about half a player, which
keeps the corrected error at 3--4\% over a whole match, provided close-ups are filtered out, which a simple filter
did on eight full broadcasts to within 1 percentage point. Team-wide pressing metrics from broadcast data should be
corrected, or compared only between broadcasts with similar framing.

\section*{Acknowledgements}
We thank IDSSE (Bassek et al.), Metrica Sports, SkillCorner, the SoccerNet team and the SportsMOT authors for making
their data openly available.

\section*{Author contributions (CRediT)}
\ifanonymous
Removed for double-anonymous review; provided in the submission system.
\else
Varun Inamdar: Conceptualization, Methodology, Software, Validation, Formal analysis, Investigation, Data curation,
Visualization, Writing -- original draft, Writing -- review \& editing.
\fi

\section*{Ethics approval and consent}
Not applicable. The study used only publicly available secondary data released by their providers for research
(match tracking and broadcast labels of professional football). No participants were recruited, and no personal data
beyond those published by the providers were processed.

\section*{Pre-registration and reporting guideline}
The study was not pre-registered. No EQUATOR reporting guideline applies to this simulation-based measurement design;
instead, every analysis setting, the code and the statistical output (Appendix~\ref{app:stats}) are reported in full.

\section*{Disclosure statement}
The authors report there are no competing interests to declare.

\section*{Funding}
No funding was obtained for this study.

\section*{Declaration of generative AI use}
Generative AI (Claude, Anthropic) was used as a coding and writing assistant: to write and test analysis code under
the author's direction, and to draft and edit parts of the manuscript text. The author specified the analyses, checked
the code, results and text, and takes full responsibility for the content of this article. No data were generated by
AI, and all reported numbers come from the archived code and its outputs.

\section*{Data availability statement}
All data are available from their original sources and are not redistributed: IDSSE \citep{bassek2025idsse}, Metrica
Sports \citep{metrica_sample}, SkillCorner \citep{skillcorner_opendata}, SoccerNet-GSR \citep{somers2024gsr},
SoccerNet-v2 \citep{deliege2021soccernetv2} and SportsMOT \citep{cui2023sportsmot}. The full SoccerNet broadcast
videos are available from the SoccerNet team under a non-disclosure agreement; only aggregate statistics derived from
them are reported.

\section*{Code availability statement}
\ifanonymous
The code, the fitted camera parameters and the scripts that rebuild every table and figure are openly available under
the Apache-2.0 licence and archived with a DOI. The link is withheld for double-anonymous review and will be added on
acceptance.
\else
The code, the fitted camera parameters and the scripts that rebuild every table and figure are openly available under
the Apache-2.0 licence \citep{pitchlens}, archived on Zenodo (version 1.0.0,
\url{https://doi.org/10.5281/zenodo.24557592}). The repository's README maps every result to its command.
\fi

\bibliographystyle{apacite}
\bibliography{refs}

\appendix
\section{Statistical output}
\label{app:stats}

\begin{table}[ht]
\centering
\caption{Fitted models. Intervals are 95\% percentile bootstrap intervals; the resampling unit is given per row.}
\label{tab:models}
\small
\begin{tabular}{p{3.4cm}p{3.9cm}p{1.9cm}p{2.6cm}}
\toprule
Model & Estimate [95\% CI] & Fit & Data, resampling \\
\midrule
Camera pan: $\phi = g_b\theta^{\text{ball}} + g_p\bar\theta^{\text{players}} + o$ & $g_b = 0.57$, $g_p = 0.56$ &
$R^2 = 0.97$ & 10 SkillCorner matches \\
Correction: bias (\%) on visible players & slope 4.12 [3.68, 4.52], intercept $-79.3$ [$-85.4$, $-72.9$] &
$R^2 = 0.76$, residual SD 2.7 pp & 400 team-broadcasts; cameras and matches \\
Correction, IDSSE / Metrica only & slope 3.9 / 4.2 & $r = 0.98$ / 0.98 & 20 cameras per dataset \\
Carrier-pressure bias (\%) on camera lag (s) & slope $-10.2$ [$-11.7$, $-9.1$] per s, intercept $-0.9$ & $r = -0.97$
& 20 cameras (mean of both datasets); cameras \\
Team effect: bias on visible players + [Fortuna] & Fortuna $-1.9$ pp [$-5.3$, $+1.5$] & -- & 400 team-broadcasts;
matches \\
Count calibration: labels on count & $a = 0.943$, $b = 1.52$ & -- & 29 SoccerNet-GSR tuning clips \\
\bottomrule
\end{tabular}
\end{table}

\begin{table}[ht]
\centering
\caption{Bias of active pressers per pressing team: mean and range over the 20 stadium cameras, and the mean absolute
error after correction with the camera and match held out (all values in \%).}
\label{tab:teams}
\small
\resizebox{\textwidth}{!}{%
\begin{tabular}{lcccc}
\toprule
Pressing team & Matches & Mean bias & Range over cameras & Corrected MAE \\
\midrule
SSV Jahn Regensburg & 1 & $-27.9$ & $-34.7$ to $-17.6$ & 5.9 \\
FC St.\ Pauli & 1 & $-24.6$ & $-30.3$ to $-15.5$ & 2.3 \\
Fortuna D\"usseldorf & 5 & $-23.9$ & $-34.3$ to $-9.7$ & 3.6 \\
Metrica game 3, team A & 1 & $-23.5$ & $-30.1$ to $-13.5$ & 3.6 \\
1.\ FC N\"urnberg & 1 & $-22.8$ & $-28.8$ to $-12.8$ & 2.1 \\
1.\ FC Kaiserslautern & 1 & $-22.7$ & $-29.7$ to $-12.8$ & 1.7 \\
Bayer 04 Leverkusen & 1 & $-21.9$ & $-28.5$ to $-11.9$ & 1.6 \\
Metrica game 3, team B & 1 & $-21.6$ & $-28.4$ to $-11.0$ & 1.5 \\
F.C.\ Hansa Rostock & 1 & $-20.6$ & $-26.7$ to $-11.3$ & 1.3 \\
1.\ FC K\"oln & 1 & $-19.9$ & $-25.3$ to $-11.8$ & 3.5 \\
Metrica games 1--2, home & 2 & $-19.5$ & $-28.0$ to $-8.6$ & 1.2 \\
VfL Bochum 1848 & 1 & $-18.5$ & $-23.3$ to $-9.9$ & 5.8 \\
Metrica games 1--2, away & 2 & $-18.4$ & $-25.5$ to $-8.2$ & 2.0 \\
FC Bayern M\"unchen & 1 & $-17.8$ & $-23.4$ to $-9.8$ & 6.3 \\
\bottomrule
\end{tabular}}
\end{table}

\begin{table}[ht]
\centering
\caption{Main-camera filter per broadcast. Default filter: the filter of the Methods (grass share at least 0.2,
median person at most 300\,px, whole shots). Broadcast filter: thresholds tuned on the other seven broadcasts and
scored on this one (held out). Shift: change of the corrected pressing value in percentage points. EPL: Premier
League; UCL: Champions League.}
\label{tab:filter-games}
\small
\resizebox{\textwidth}{!}{%
\begin{tabular}{llcccc}
\toprule
 & & Default filter & \multicolumn{3}{c}{Broadcast filter (held out)} \\
\cmidrule(lr){3-3}\cmidrule(lr){4-6}
League & Match & Shift & Live main removed & Non-main kept & Shift \\
\midrule
EPL & Manchester United 4--1 Leicester & $-6.5$ & 1.8\% & 12.0\% & $+0.07$ \\
EPL & Swansea 1--2 Liverpool & $-5.6$ & 0.9\% & 11.4\% & $+0.60$ \\
UCL & Manchester City 3--1 Barcelona & $-11.6$ & 1.8\% & 11.6\% & $-0.27$ \\
Ligue 1 & Paris SG 3--0 Metz & $-9.3$ & 2.8\% & 18.0\% & $-0.42$ \\
Bundesliga & Bayern Munich 5--1 Wolfsburg & $-6.7$ & 1.3\% & 7.8\% & $+0.19$ \\
Serie A & AS Roma 2--0 Empoli & $-8.7$ & 2.2\% & 8.5\% & $+0.43$ \\
La Liga & Real Madrid 3--2 Valencia & $-3.7$ & 1.6\% & 4.9\% & $-0.27$ \\
La Liga & Dep.\ La Coru\~na 0--2 Real Madrid & $-3.0$ & 1.3\% & 5.0\% & $+0.13$ \\
\bottomrule
\end{tabular}}
\end{table}

\end{document}
